\documentclass[12pt]{article}

% =========================
% PREAMBULO (el mismo de la tesis)
% =========================

% Idioma y codificación
\usepackage[spanish,es-tabla]{babel}
\usepackage[utf8]{inputenc}
\usepackage[T1]{fontenc}
\usepackage{mathptmx}
\usepackage{amsmath}
\usepackage{amssymb}
% Tablas bonitas
\usepackage{booktabs}
\usepackage{makecell}
\usepackage{tabularx}
\usepackage[table]{xcolor}
\usepackage{longtable}
\usepackage{array}
\usepackage{amsmath}
\definecolor{headerblue}{RGB}{26,97,133}
\definecolor{lineblue}{RGB}{82,172,220}
\definecolor{rowblue}{RGB}{173,211,230}

\arrayrulecolor{lineblue}
\setlength{\arrayrulewidth}{0.5pt}

\usepackage[ruled,vlined]{algorithm2e}
\newcommand{\headercell}[1]{\textcolor{white}{\textbf{#1}}}
% Control de figuras/tablas y referencias
\usepackage{caption}
\usepackage{subcaption} % subfiguras
\usepackage{tikz}
\usetikzlibrary{arrows.meta, positioning}
% Para que [H] funcione y no se muevan tanto
\usepackage{float}
\usepackage{tikz}
\usepackage{xcolor}
% Links (índice clickeable, referencias clickeables)
\usepackage[colorlinks=true,linkcolor=black,citecolor=black,urlcolor=blue]{hyperref}

% Listas mejoradas
\usepackage{enumitem}
\definecolor{azulESCOM}{RGB}{82,172,220}
% Código (recomendado)
\usepackage{listings}
\usepackage{xcolor}
% Bibliografía IEEE

\usepackage[backend=biber,style=ieee]{biblatex}

% Evita los parentesis vacios "()" en referencias web sin fecha de publicacion
\usepackage{xpatch}
\xpatchbibdriver{online}
  {\printtext[parens]{\usebibmacro{date}}}
  {\iffieldundef{year}{}{\printtext[parens]{\usebibmacro{date}}}}
  {}{}
\addbibresource{references.bib}
\setlength{\parindent}{0pt}%sangria
\setlength{\parskip}{0.4em}

\usepackage{csquotes}

% Gráficos y márgenes
\usepackage{graphicx}
\usepackage[
a4paper,
left=3cm,
right=2cm,
top=2cm,
bottom=2cm
]{geometry}

% Numeracion de tablas y figuras por seccion (Tabla 3.1, Figura 3.2, ...)
\numberwithin{table}{section}
\numberwithin{figure}{section}
\numberwithin{equation}{section}

% Margen de holgura para lineas dificiles de partir (URLs largas en la bibliografia)
\setlength{\emergencystretch}{3em}


% ---- Ajustes propios del manual ----
% Marca visible de "texto por redactar" (se ve en el PDF en gris cursiva)
\newcommand{\pendiente}[1]{\textcolor{gray}{\textit{[Por redactar: #1]}}}
% Estilo para comandos de consola y archivos de configuracion
\lstdefinestyle{consola}{
  basicstyle=\ttfamily\small,
  backgroundcolor=\color{gray!10},
  frame=single, rulecolor=\color{lineblue},
  breaklines=true, columns=fullflexible,
  showstringspaces=false, keepspaces=true
}
\lstset{style=consola}

\begin{document}

% =========================
% PRELIMINARES
% =========================
%% Portadas (institucional y con resumen) -- las mismas de la tesis

% =========================
% PORTADA
% =========================


\newgeometry{top=2cm, bottom=2cm, right=2cm, left=3cm}

\begin{titlepage}
	\thispagestyle{empty}
	\centering
	
	% --- LÍNEAS Y LOGOS ---
	\begin{tikzpicture}[remember picture, overlay]
		
		% Logo IPN alineado con las líneas azules
		\node[anchor=north] 
		at ([xshift=1.60cm, yshift=-1.25cm]current page.north west) {
			\includegraphics[width=2.25cm]{images/ipn.png}
		};
		
		% Tres líneas azules entre los logos, como en la plantilla
		% No atraviesan las imágenes
		\foreach \x in {1.45,1.60,1.75} {
			\draw[azulESCOM, line width=1.5pt] 
			([xshift=\x cm, yshift=-3.85cm]current page.north west) -- 
			([xshift=\x cm, yshift=4.20cm]current page.south west);
		}
		
		% Logo ESCOM alineado con las líneas azules
		\node[anchor=south] 
		at ([xshift=1.60cm, yshift=2.0cm]current page.south west) {
			\includegraphics[width=2.45cm]{images/ESCOM.png}
		};
		
	\end{tikzpicture}
	
	\vspace*{0.85cm}
	
	% --- BLOQUE INSTITUCIONAL ---
	{\fontsize{20}{24}\selectfont \textbf{INSTITUTO POLITÉCNICO NACIONAL}}\\[0.35cm]
	{\fontsize{16}{19.2}\selectfont \textbf{ESCUELA SUPERIOR DE CÓMPUTO}}\\[0.55cm]
	{\fontsize{20}{24}\selectfont \textbf{ESCOM}}\\[1.45cm]
	
	% --- BLOQUE DEL PROYECTO ---
	{\fontsize{14}{16.8}\selectfont \textit{Trabajo Terminal}}\\[0.75cm]
	
	{\fontsize{14}{16.8}\selectfont 2026-B009}\\[0.75cm]
	
	{\fontsize{16}{19.2}\selectfont \bfseries
		``Manual de Instalación de la Plataforma Web para la Denuncia Segura y Prevención de Riesgos Digitales''
	}\\[0.6cm]
	
	\vfill
	
	% --- BLOQUE DE AUTOR ---
	{\fontsize{14}{16.8}\selectfont \textit{Presentan}}\\[0.4cm]
	{\fontsize{14}{16.8}\selectfont \textbf{Edgar Jair Flores Espinosa}}\\[.5cm]
	{\fontsize{14}{16.8}\selectfont \textbf{Jose Bryan Omar Jimenez Velazquez}}\\[.5cm]
	{\fontsize{14}{16.8}\selectfont \textbf{Leilani Shareni Santos Rivera}}\\[.5cm]

	
	% --- BLOQUE DE DIRECTOR ---
	{\fontsize{14}{16.8}\selectfont \textit{Directores}}\\[0.4cm]
	{\fontsize{14}{16.8}\selectfont \textbf{Dr. Jessie Paulina Guzman Flores}}\\[.3cm]
	{\fontsize{14}{16.8}\selectfont \textbf{M. en C. Ulises Velez Saldaña}}\\[.3cm]
	\vfill
	
	% --- PIE DE PÁGINA ---
	{\fontsize{12}{14.4}\selectfont 29 de Diciembre de 2026}
	
\end{titlepage}

\restoregeometry






	
	% =========================
	% PORTADA ADICIONAL CON RESUMEN
	% =========================
	\begin{titlepage}
		\thispagestyle{empty}
		
		\begin{center}
			
			% Encabezado bien alineado
			\noindent
			\begin{minipage}[c]{0.18\textwidth}
				\centering
				\includegraphics[width=2.4cm]{images/ipn.png}
			\end{minipage}
			\begin{minipage}[c]{0.62\textwidth}
				\centering
				{\fontsize{20}{22}\selectfont \textbf{INSTITUTO POLITÉCNICO NACIONAL}\par}
				{\fontsize{16}{18}\selectfont \textbf{ESCUELA SUPERIOR DE CÓMPUTO}\par}
				{\fontsize{16}{18}\selectfont \textbf{SUBDIRECCIÓN ACADÉMICA}\par}
			\end{minipage}
			\begin{minipage}[c]{0.18\textwidth}
				\centering
				\includegraphics[width=3.0cm]{images/ESCOM.png}
			\end{minipage}
			
			\vspace{0.9cm}
			
			\noindent
			\begin{minipage}{0.45\textwidth}
				\raggedright
				{\fontsize{14}{16}\selectfont No. de TT: 2026-B009}
			\end{minipage}
			\begin{minipage}{0.50\textwidth}
				\raggedleft
				{\fontsize{14}{16}\selectfont Fecha de presentación de TT-II: 05 de junio de 2026}
			\end{minipage}
			
			\vspace{0.55cm}
			
			{\fontsize{14}{16}\selectfont Documento Técnico\par}
			
			\vspace{0.65cm}
			
			\begin{minipage}{0.88\textwidth}
				\centering
				{\fontsize{16}{18}\selectfont \textbf{``Manual de Instalación de la iPlataforma Web para la Denuncia Segura y Prevención de Riesgos Digitales''}\par}
			\end{minipage}
			
			\vspace{1.25cm}
			
			{\fontsize{14}{16}\selectfont \textit{Presentan:}\par}
			\vspace{0.15cm}
			{\fontsize{14}{16}\selectfont \textbf{Edgar Jair Flores Espinosa, Jose Bryan Omar Jimenez Velazquez, Leilani Shareni Santos Rivera}\par}
			
			\vspace{1.35cm}
			
			{\fontsize{14}{16}\selectfont \textit{Directores:}\par}
			
			\vspace{0.45cm}
			
			{\fontsize{14}{16}\selectfont \textbf{Dr. Jessie Paulina Guzman Flores}\par}
			{\fontsize{14}{16}\selectfont \textbf{M. en C. Ulises Velez Saldaña}\par}
						
			\vspace{1.1cm}
			
			{\fontsize{14}{16}\selectfont \textbf{MANUAL DE INSTALACION}\par}
			
		\end{center}
		
		\vspace{0.25cm}
		
		\vfill
		
		{\fontsize{12}{14}\selectfont
			\rule{6cm}{0.4pt}\\[-0.05cm]
			\textsuperscript{1}\texttt{josebryanomar2004@gmail.com}
		}
		
	\end{titlepage}


%% Indice general, de figuras y de tablas
	% Numeración romana para preliminares
	
	\pagenumbering{roman}
	
	\tableofcontents
	\newpage
	
	\listoffigures
	\newpage
	
	\listoftables
	\newpage
	
	% Numeración arábiga para el contenido
	\pagenumbering{arabic}
	
	


% =========================
% CONTENIDO DEL MANUAL
% =========================
\section{Introducción}

\subsection{Objetivo del manual}
\pendiente{qué se logra al terminar el manual: la plataforma completa funcionando en https://defensoria-escom.ddns.net}

\subsection{Alcance}
\pendiente{qué cubre (BD, microservicios, frontends, router, HTTPS, modelo de IA) y qué no cubre (carga del histórico, que la realiza otra persona; servicios de primer contacto y subdefensoría, que pertenecen a otros equipos)}

\subsection{A quién va dirigido y conocimientos previos}
\pendiente{administrador de sistemas con nociones de Linux, SSH, contenedores y redes}

\subsection{Convenciones del documento}
\pendiente{cómo se muestran los comandos (\texttt{\#} = root, \texttt{\$} = usuario), rutas, valores que se deben sustituir como \texttt{<IP\_BD>}; las contraseñas nunca se escriben en este documento}

\subsection{Arquitectura general del despliegue}
\pendiente{diagrama TikZ de las 3 VPS: usuario $\rightarrow$ router-nginx (Frontend) $\rightarrow$ microservicios (Backend) $\rightarrow$ PostgreSQL (BD)}

\begin{table}[H]
\centering
\caption{Servidores que componen el despliegue}
\begin{tabular}{lll}
\toprule
\textbf{Rol} & \textbf{Dirección IP} & \textbf{Contenido} \\
\midrule
Base de datos & 169.58.62.99 & defensoria-db (5432), historico-db (5433) \\
Backend & 156.67.26.73 & Microservicios Spring Boot (8082--8094) \\
Frontend & 169.58.62.111 & router-nginx (80/443) y frontends Angular \\
\bottomrule
\end{tabular}
\end{table}

\subsection{Flujo de una petición}
\pendiente{ejemplo: el quejoso envía una queja $\rightarrow$ nginx /api/quejas/ $\rightarrow$ queja-service:8084 $\rightarrow$ defensoria-db; notificaciones por correo}

\subsection{Orden recomendado de instalación}
\begin{enumerate}
  \item Preparar las tres VPS (sección \ref{sec:preparacion}).
  \item Servidor de base de datos.
  \item Compilar los artefactos en el equipo de desarrollo.
  \item Servidor backend y modelo de IA.
  \item Servidor frontend, dominio y HTTPS.
  \item Datos iniciales, verificación y puesta en operación.
\end{enumerate}

\section{Requisitos y versiones}

\subsection{Recursos de cada VPS}
\pendiente{CPU, RAM, disco y sistema operativo de cada servidor (Contabo, Ubuntu, kernel 6.8); salida de \texttt{lscpu}, \texttt{free -h}, \texttt{df -h}}

\subsection{Software base en los servidores}
\begin{table}[H]
\centering
\caption{Software base}
\begin{tabular}{lll}
\toprule
\textbf{Componente} & \textbf{Versión} & \textbf{Dónde} \\
\midrule
Ubuntu Server & \pendiente{versión} & Las 3 VPS \\
Podman & \pendiente{\texttt{podman --version}} & Las 3 VPS \\
iptables & \pendiente{versión} & BD y Backend \\
Certbot & \pendiente{versión} & Frontend \\
\bottomrule
\end{tabular}
\end{table}

\subsection{Imágenes de contenedor}
\begin{table}[H]
\centering
\caption{Imágenes base utilizadas}
\begin{tabular}{lll}
\toprule
\textbf{Imagen} & \textbf{Versión} & \textbf{Uso} \\
\midrule
postgres & 16-alpine (16.15) & defensoria-db, historico-db \\
eclipse-temurin & 21-jre-jammy & Todos los microservicios \\
nginx & alpine (1.31.x) & router-nginx y frontends \\
\bottomrule
\end{tabular}
\end{table}

\subsection{Herramientas de desarrollo (compilación)}
\begin{table}[H]
\centering
\caption{Versiones de compilación}
\begin{tabular}{ll}
\toprule
\textbf{Herramienta / librería} & \textbf{Versión} \\
\midrule
Java (JDK) & 21 \\
Maven & \pendiente{versión} \\
Spring Boot & 3.5.16 (primercontacto y subdefensoria: 3.3.5) \\
springdoc-openapi & 2.8.5 \\
jjwt & 0.12.6 \\
Node.js / npm & \pendiente{versión} \\
Angular / Angular CLI & 21.2 / 21.2.8 \\
\bottomrule
\end{tabular}
\end{table}

\subsection{Puertos utilizados}
\pendiente{referencia a la tabla completa de puertos del Anexo \ref{sec:anexos}}

\subsection{Cuentas y servicios externos}
\pendiente{dominio No-IP, cuenta de correo SMTP para notificaciones, acceso SSH root a las 3 VPS}

\section{Preparación de las VPS}\label{sec:preparacion}

\subsection{Primer acceso por SSH}
\pendiente{conexión, cambio de contraseña root, recomendación de llaves SSH}

\subsection{Actualización del sistema y zona horaria}
\begin{lstlisting}
# apt update && apt upgrade -y
# timedatectl set-timezone America/Mexico_City
\end{lstlisting}

\subsection{Nombre de host}
\pendiente{vmi3633025 (BD), vmi3633021 (Backend), vmi3633026 (Frontend)}

\subsection{Instalación de Podman}
\pendiente{\texttt{apt install podman}, verificación con \texttt{podman info}, registro por defecto docker.io}

\subsection{Habilitar el reinicio automático de contenedores}
\pendiente{\texttt{podman-restart.service} y la política \texttt{--restart always}}

\subsection{Estructura de directorios en el servidor}
\pendiente{carpeta de trabajo de cada VPS (scripts, \texttt{config-files/}, \texttt{jars/}, \texttt{database-scripts/}, respaldos)}

\subsection{Transferencia de archivos al servidor}
\pendiente{uso de \texttt{scp} / WinSCP / FileZilla desde la carpeta \texttt{Produccion}}

\subsection{Consideraciones de firewall}
\pendiente{por qué ufw no basta con Podman (el DNAT lo omite) y por qué se usan reglas en la tabla raw de iptables limitadas a la interfaz pública}

\section{Servidor de base de datos (169.58.62.99)}

\subsection{Archivos necesarios}
\pendiente{\texttt{podman-compose-bd.sh}, \texttt{database-scripts/init/00-init.sql} y \texttt{database-scripts/seeds/}}

\subsection{Descripción del script podman-compose-bd.sh}
\pendiente{comandos disponibles: up, reset, seed, firewall, firewall-boot, backup, restore, psql-*}

\subsection{Volúmenes y persistencia de datos}
\pendiente{dónde viven los datos de PostgreSQL y qué pasa al recrear el contenedor}

\subsection{Creación de defensoria-db (puerto 5432)}
\pendiente{comando, variables POSTGRES\_*, base \texttt{defensoria\_db}}

\subsection{Creación de historico-db (puerto 5433)}
\pendiente{base \texttt{defensoria\_historico\_db}; la carga de datos históricos queda a cargo de otra persona}

\subsection{Script de inicialización y esquema}
\pendiente{qué crea 00-init.sql y qué tablas genera Hibernate (\texttt{ddl-auto: update}) al arrancar los servicios}

\subsection{Carga de catálogos semilla}
\pendiente{01-dependencias, 02-chatbot, 03-personal\_test (este último solo para pruebas)}

\subsection{Firewall: acceso solo desde el Backend}
\pendiente{\texttt{./podman-compose-bd.sh firewall} y \texttt{firewall-boot}; los puertos 5432/5433 solo aceptan la IP 156.67.26.73}

\subsection{Verificación}
\begin{lstlisting}
# podman ps
# ./podman-compose-bd.sh psql-defensoria   # \dt
\end{lstlisting}

\section{Compilación de los artefactos}

\subsection{Preparación del equipo de desarrollo}
\pendiente{JDK 21, Maven, Node.js y Angular CLI; variables JAVA\_HOME}

\subsection{Estructura del repositorio (carpeta Produccion)}
\pendiente{Backend/, Frontend*/, Modelo-Java/, nginx/, docs/}

\subsection{Compilación de los microservicios con Maven}
\begin{lstlisting}
$ cd Backend/queja-service
$ mvn clean package -DskipTests
$ ls target/*.jar
\end{lstlisting}
\pendiente{repetir para cada servicio: auth, queja, notificaciones, catalogo, admin, revision, chatbot, denunciado}

\subsection{Compilación de los frontends Angular}
\begin{lstlisting}
$ cd Frontend
$ npm ci
$ npx ng build --configuration production
\end{lstlisting}
\pendiente{Frontend (defensoria-web), Frontend-Revision (/revision/), Frontend-Admin (/admin/); el \texttt{base-href} de cada uno}

\subsection{Artefactos resultantes y a dónde se suben}
\pendiente{tabla: artefacto $\rightarrow$ servidor $\rightarrow$ carpeta destino}

\section{Servidor backend (156.67.26.73)}

\subsection{Archivos necesarios}
\pendiente{jars, Dockerfiles, \texttt{config-files/}, \texttt{podman-compose.sh}, \texttt{firewall-backend.sh}}

\subsection{Archivos de configuración externos (config-files)}
\pendiente{un YAML por servicio montado en \texttt{/config}; selección con \texttt{SPRING\_CONFIG\_NAME}}

\subsubsection{Conexión a la base de datos}
\pendiente{URL jdbc hacia 169.58.62.99:5432 y :5433}
\subsubsection{Secreto JWT compartido}
\pendiente{el mismo \texttt{jwt.secret} en todos los servicios}
\subsubsection{Correo SMTP y tiempos de espera}
\pendiente{auth y notificaciones; timeouts de correo y de Feign}
\subsubsection{URLs entre microservicios}
\pendiente{cómo se llaman entre sí (puertos publicados)}

\subsection{Dockerfile de los microservicios}
\pendiente{base eclipse-temurin:21-jre-jammy; caso especial admin-service con postgresql-client-16 (PGDG) para respaldos}

\subsection{El script podman-compose.sh}
\pendiente{construcción de imágenes, levantar/detener servicios, logs}

\subsection{Límites de memoria y JVM}
\begin{table}[H]
\centering
\caption{Memoria asignada por microservicio}
\begin{tabular}{lll}
\toprule
\textbf{Servicio} & \textbf{Límite contenedor} & \textbf{Heap JVM} \\
\midrule
queja-service & 1 GB & -Xmx640m \\
denunciado-service & 1 GB & \pendiente{} \\
auth, notificaciones, catalogo, admin, revision, chatbot & 512 MB & -Xmx320m \\
primercontacto, subdefensoria & 250 MB & \pendiente{} \\
antecedentes-service (IA) & 250 MB & -Xmx150m \\
\bottomrule
\end{tabular}
\end{table}
\pendiente{uso de \texttt{-XX:+ExitOnOutOfMemoryError} y \texttt{--restart always}}

\subsection{Despliegue de cada microservicio}
\pendiente{una subsección corta por servicio: propósito, puerto, dependencias, comando y verificación}
\subsubsection{auth-service (8083)}
\subsubsection{queja-service (8084)}
\subsubsection{notificaciones-service (8085)}
\subsubsection{catalogo-service (8086)}
\subsubsection{admin-service (8087)}
\subsubsection{revision-service (8088)}
\subsubsection{chatbot-service (8089)}
\subsubsection{denunciado-service (8094)}
\subsubsection{historico-service (8092)}
\pendiente{no se levanta por defecto}
\subsubsection{Servicios de otros equipos: primercontacto (8082) y subdefensoria (8091)}

\subsection{Firewall del backend}
\pendiente{\texttt{firewall-backend.sh aplicar|boot|quitar}; rango 8082:8094 solo desde el Frontend y solo en la interfaz pública para no bloquear el tráfico entre contenedores}

\subsection{Verificación}
\pendiente{\texttt{podman ps}, \texttt{podman stats}, \texttt{curl localhost:<puerto>/actuator/health} o Swagger}

\section{Modelo de IA: antecedentes-service (8093)}\label{sec:modelo}

\subsection{Descripción del modelo}
El microservicio \texttt{antecedentes-service} contiene el modelo de inteligencia artificial de la plataforma. No depende de servicios externos ni de GPU: se ejecuta dentro de la JVM con dos técnicas de procesamiento de lenguaje natural.
\begin{itemize}
  \item \textbf{Búsqueda de antecedentes:} representa cada queja con TF-IDF y la compara con el corpus histórico mediante similitud coseno. Devuelve las quejas más parecidas por encima de un umbral.
  \item \textbf{Resumen extractivo:} aplica TextRank para elegir las oraciones más representativas de una queja larga, descartando las redundantes.
  \item \textbf{Aprendizaje incremental:} las quejas nuevas se agregan al corpus y se guardan en \texttt{quejas\_aprendidas.json}, sin reentrenar ni reiniciar.
\end{itemize}
\pendiente{diagrama del flujo: texto $\rightarrow$ normalización $\rightarrow$ vector TF-IDF $\rightarrow$ coseno $\rightarrow$ top-k}

\subsection{Quién lo consume}
\pendiente{revision-service lo llama durante el turnado; Frontend-Revision muestra los antecedentes y el resumen en la pantalla de turnado}

\subsection{Estructura del proyecto Modelo-Java}
\begin{table}[H]
\centering
\caption{Paquetes del modelo}
\begin{tabular}{ll}
\toprule
\textbf{Paquete} & \textbf{Contenido} \\
\midrule
\texttt{antecedentes\_service} & API REST, configuración (\texttt{IaProperties}) y servicio \\
\texttt{ia.dataset} & Lectura de los JSON del corpus (\texttt{LectorDataset}) \\
\texttt{ia.normalizacion} & Limpieza y normalización del texto \\
\texttt{ia.tfidf} & Índice TF-IDF y vectores dispersos \\
\texttt{ia.antecedentes} & Buscador por similitud coseno \\
\texttt{ia.resumen} & Resumen extractivo TextRank \\
\bottomrule
\end{tabular}
\end{table}

\subsection{Conjuntos de datos (Dataset)}
\begin{table}[H]
\centering
\caption{Archivos del corpus}
\begin{tabular}{lp{8.5cm}}
\toprule
\textbf{Archivo} & \textbf{Uso} \\
\midrule
\texttt{quejas\_sinteticas.json} & Corpus histórico base del índice \\
\texttt{quejas\_prueba\_30.json} & Quejas nuevas de prueba \\
\texttt{quejas\_aprendidas.json} & Lo escribe el propio servicio al aprender quejas nuevas (no se debe sobrescribir en producción) \\
\bottomrule
\end{tabular}
\end{table}
\pendiente{cómo se generaron los datos sintéticos con \texttt{generar\_datasets.py} y \texttt{generar\_quejas\_variadas.py}}

\subsection{Compilación}
\begin{lstlisting}
$ cd Modelo-Java
$ mvn clean package -DskipTests
$ cp target/antecedentes-service-0.1.0.jar antecedentes-service.jar
\end{lstlisting}
Spring Boot 3.5.16, Java 21, solo \texttt{spring-boot-starter-web}.

\subsection{Archivos en el servidor backend}
\begin{table}[H]
\centering
\caption{Ubicación de los archivos del modelo (base: /apps/aplicaciones/defensoria/back)}
\begin{tabular}{ll}
\toprule
\textbf{Archivo} & \textbf{Ruta} \\
\midrule
JAR & \texttt{artifact/antecedentes-service.jar} \\
Configuración & \texttt{config-files/antecedentes-service/config/antecedentes-service.yml} \\
Datasets & \texttt{modelo/Dataset/} \\
Script & \texttt{podman-compose-modelo.sh} \\
\bottomrule
\end{tabular}
\end{table}

\subsection{Parámetros de configuración}
\begin{table}[H]
\centering
\caption{Parámetros del modelo (prefijo \texttt{ia.})}
\begin{tabular}{llp{6.5cm}}
\toprule
\textbf{Parámetro} & \textbf{Valor} & \textbf{Significado} \\
\midrule
\texttt{busqueda.umbral} & 0.25 & Similitud mínima para considerar un antecedente (subir a 0.70 al comparar quejas completas con un corpus grande) \\
\texttt{busqueda.top-k} & 5 & Antecedentes devueltos por defecto \\
\texttt{busqueda.top-k-maximo} & 50 & Límite que puede pedir el cliente \\
\texttt{resumen.numero-oraciones} & 3 & Oraciones del resumen por defecto \\
\texttt{resumen.numero-oraciones-maximo} & 20 & Límite de oraciones \\
\texttt{resumen.redundancia} & 0.60 & Similitud a partir de la cual una oración se descarta por repetida \\
\texttt{dataset.*} & \texttt{/app/data/...} & Rutas de los JSON dentro del contenedor \\
\bottomrule
\end{tabular}
\end{table}

\subsection{Despliegue con podman-compose-modelo.sh}
\begin{lstlisting}
# sudo bash podman-compose-modelo.sh up       # construye y levanta
# sudo bash podman-compose-modelo.sh status   # estado + health
# sudo bash podman-compose-modelo.sh logs
# sudo bash podman-compose-modelo.sh restart | delete
\end{lstlisting}
El contenedor publica \texttt{8093:8080}, se limita a 250\,MB con \texttt{-Xmx150m -Xms64m -XX:+UseSerialGC} y monta la configuración en \texttt{/app/config} y los datasets en \texttt{/app/data}.

\subsection{Endpoints}
\begin{table}[H]
\centering
\caption{API del modelo (\texttt{/api/antecedentes})}
\begin{tabular}{llp{7cm}}
\toprule
\textbf{Método} & \textbf{Ruta} & \textbf{Función} \\
\midrule
POST & \texttt{/buscar} & Antecedentes similares a un texto (\texttt{texto}, \texttt{umbral}, \texttt{topK}) \\
POST & \texttt{/resumir} & Resumen extractivo de una queja \\
POST & \texttt{/entrenar} & Agrega una queja al corpus \\
POST & \texttt{/entrenar-batch} & Agrega varias quejas \\
GET & \texttt{/corpus} & Tamaño y estado del corpus \\
GET & \texttt{/health} & Verificación de funcionamiento \\
\bottomrule
\end{tabular}
\end{table}

\subsection{Verificación}
\begin{lstlisting}
# curl -s http://127.0.0.1:8093/api/antecedentes/health
# curl -s -X POST http://127.0.0.1:8093/api/antecedentes/buscar \
    -H "Content-Type: application/json" -d @pruebas/buscar.json
\end{lstlisting}
\pendiente{respuesta esperada; la carpeta \texttt{Modelo-Java/pruebas} trae más ejemplos (resumir, entrenar)}

\subsection{Respaldo del aprendizaje}
\pendiente{respaldar \texttt{modelo/Dataset/quejas\_aprendidas.json} junto con la base de datos}

\section{Servidor frontend (169.58.62.111)}

\subsection{Archivos necesarios}
\pendiente{carpetas \texttt{dist/} de cada frontend, Dockerfiles nginx, carpeta \texttt{nginx/config}}

\subsection{Contenedores de los frontends}
\begin{table}[H]
\centering
\caption{Contenedores web}
\begin{tabular}{lll}
\toprule
\textbf{Contenedor} & \textbf{Puerto} & \textbf{Ruta pública} \\
\midrule
defensoria-web & 8090 & / (portal del quejoso y del denunciado) \\
revision-web & 8092 $\rightarrow$ 80 & /revision/ (login unificado del personal) \\
admin-web & 8091 $\rightarrow$ 80 & /admin/ (Panel de la Defensora) \\
\bottomrule
\end{tabular}
\end{table}

\subsection{Construcción y arranque de cada frontend}
\pendiente{\texttt{podman build} y \texttt{podman run} de cada uno}

\subsection{router-nginx}
\pendiente{contenedor con \texttt{--network host} que atiende 80/443}

\subsection{Archivo defensoria.conf explicado}
\subsubsection{Bloque del puerto 80 y redirección a HTTPS}
\subsubsection{Bloque 443 y certificados}
\subsubsection{Rutas hacia los frontends}
\subsubsection{Rutas /api/ hacia los microservicios}
\pendiente{tabla ruta $\rightarrow$ servicio:puerto (ver Anexo)}
\subsubsection{Rutas bloqueadas (403)}
\pendiente{\texttt{/api/notificaciones/enviar} y \texttt{/registrar} solo para uso interno}
\subsubsection{Tamaño máximo de carga y tiempos de espera}
\pendiente{\texttt{client\_max\_body\_size} para evidencias, \texttt{proxy\_read\_timeout}, cabecera X-Real-IP}

\subsection{Recarga de la configuración}
\begin{lstlisting}
# podman exec router-nginx nginx -t
# podman exec router-nginx nginx -s reload
\end{lstlisting}

\section{Dominio y HTTPS}

\subsection{Dominio dinámico en No-IP}
\pendiente{registro de defensoria-escom.ddns.net apuntando a 169.58.62.111 y renovación del host}

\subsection{Obtención del certificado con Certbot (webroot)}
\pendiente{instalación de certbot, carpeta webroot compartida con router-nginx, comando \texttt{certonly --webroot}}

\subsection{Montaje de los certificados en router-nginx}
\pendiente{volúmenes de /etc/letsencrypt}

\subsection{Renovación automática}
\pendiente{\texttt{certbot renew --dry-run}, hook para recargar nginx}

\subsection{Verificación}
\pendiente{navegador y \texttt{curl -I https://defensoria-escom.ddns.net}}

\section{Datos iniciales y cuentas}

\subsection{Catálogos (dependencias y chatbot)}
\pendiente{cómo verificar que se cargaron}

\subsection{Roles del personal}
\pendiente{RECEPCIONISTA, DEFENSOR (Defensora), ADMIN\_SISTEMAS y demás roles de revisión; qué puede hacer cada uno}

\subsection{Cuenta de administración inicial}
\pendiente{existe una cuenta de administración que obliga a cambiar la contraseña en el primer acceso. Las credenciales se entregan por separado, no en este documento}

\subsection{Alta de personal}
\pendiente{desde el Panel de la Defensora o con SQL (INSERT en personal\_administrativo con contraseña cifrada en BCrypt)}

\subsection{Eliminar las cuentas de prueba antes de producción}
\pendiente{cuentas *.test@ipn.mx del seed 03-personal\_test}

\section{Verificación de la instalación}

\subsection{Lista de comprobación por servidor}
\pendiente{contenedores arriba, puertos escuchando, firewall activo}

\subsection{Pruebas funcionales de punta a punta}
\subsubsection{Quejoso: registrar y consultar una queja}
\subsubsection{Personal: login unificado y validación por la recepcionista}
\subsubsection{Defensora: panel, indicadores y turnado}
\subsubsection{Denunciado: registro de su respuesta}
\subsubsection{Notificaciones por correo}
\subsubsection{Chatbot}

\subsection{Pruebas de seguridad básicas}
\pendiente{puertos 5432/5433 y 8082--8094 cerrados desde Internet; rutas bloqueadas devuelven 403; solo HTTPS}

\section{Operación y mantenimiento}

\subsection{Consultar logs}
\begin{lstlisting}
# podman logs -f --tail 200 queja-service
\end{lstlisting}

\subsection{Monitoreo de recursos}
\pendiente{\texttt{podman stats}, \texttt{free -h}, \texttt{dmesg} para detectar OOM}

\subsection{Reinicio de servicios y del servidor}
\pendiente{qué se levanta solo tras reiniciar la VPS y qué hay que volver a aplicar (firewall boot)}

\subsection{Respaldos y restauración}
\subsubsection{Desde el servidor de BD (podman-compose-bd.sh backup/restore)}
\subsubsection{Desde el Panel de la Defensora (admin-service)}
\subsubsection{Política de respaldo recomendada}

\subsection{Actualizar un microservicio (redespliegue)}
\pendiente{compilar jar $\rightarrow$ subir $\rightarrow$ reconstruir imagen $\rightarrow$ recrear contenedor}

\subsection{Actualizar un frontend}
\pendiente{ng build $\rightarrow$ subir dist $\rightarrow$ reconstruir contenedor}

\subsection{Cambiar contraseñas y secretos}
\pendiente{BD, jwt.secret, SMTP: en qué archivos se cambian y qué se reinicia}

\section{Solución de problemas}

\subsection{Un microservicio se reinicia o termina con código 137 (OOM)}
\pendiente{diagnóstico con dmesg/podman stats; ajustar -m y -Xmx}

\subsection{Error 504 o peticiones que se quedan colgadas}
\pendiente{regla de firewall sin \texttt{-i <interfaz>} que bloquea tráfico entre contenedores; timeouts de RestTemplate/Feign}

\subsection{El sitio pierde HTTPS o redirige al portal principal}
\pendiente{defensoria.conf incompleto; falta una ruta location}

\subsection{No se envían correos}
\pendiente{credenciales SMTP, endpoint /api/notificaciones/enviar, timeouts}

\subsection{Error de conexión a la base de datos}
\pendiente{firewall de BD, IP permitida, contraseña en config-files}

\subsection{Fallan los respaldos desde el panel}
\pendiente{versión del cliente pg\_dump debe ser 16}

\subsection{Error 401/403 en el login del personal}
\pendiente{jwt.secret distinto entre servicios, rol sin permisos, sesión guardada en el navegador}

\section{Anexos}\label{sec:anexos}

\subsection{Tabla de puertos}
\begin{table}[H]
\centering
\caption{Puertos de la plataforma}
\begin{tabular}{llll}
\toprule
\textbf{Servicio} & \textbf{Servidor} & \textbf{Puerto} & \textbf{Ruta nginx} \\
\midrule
defensoria-db & BD & 5432 & -- \\
historico-db & BD & 5433 & -- \\
primercontacto & Backend & 8082 & \pendiente{} \\
auth-service & Backend & 8083 & /api/auth/ \\
queja-service & Backend & 8084 & /api/quejas/ \\
notificaciones-service & Backend & 8085 & /api/notificaciones/ \\
catalogo-service & Backend & 8086 & /api/catalogo/ \\
admin-service & Backend & 8087 & /api/admin/ \\
revision-service & Backend & 8088 & /api/revision/ \\
chatbot-service & Backend & 8089 & /api/chatbot/ \\
subdefensoria & Backend & 8091 & \pendiente{} \\
historico-service & Backend & 8092 & \pendiente{} \\
antecedentes (IA) & Backend & 8093 & \pendiente{} \\
denunciado-service & Backend & 8094 & /api/denunciado/ \\
defensoria-web & Frontend & 8090 & / \\
admin-web & Frontend & 8091 & /admin/ \\
revision-web & Frontend & 8092 & /revision/ \\
router-nginx & Frontend & 80, 443 & -- \\
\bottomrule
\end{tabular}
\end{table}
\pendiente{confirmar las rutas nginx contra defensoria.conf}

\subsection{Resumen de comandos por servidor}
\pendiente{chuleta de comandos más usados}

\subsection{Archivo defensoria.conf completo}
\pendiente{\texttt{\textbackslash lstinputlisting} del archivo real}

\subsection{Ejemplo de archivo de configuración de un microservicio}
\pendiente{YAML con contraseñas reemplazadas por \texttt{<CONTRASEÑA>}}

\subsection{Registro de cambios del manual}
\pendiente{versión, fecha, cambio}


% =========================
% REFERENCIAS (documentacion oficial)
% =========================
\nocite{*}
\newpage
\printbibliography[title={Referencias}]

\end{document}
